
When large language models reason step-by-step, they generate linguistic markers of uncertainty that correlate with confidence judgments. This observation motivates the present investigation of whether CoT prompting enables a form of ''self-reading'' whereby models leverage their own in-context uncertainty signals.

\subsubsection{The Calibration Problem}

LLM confidence estimates are poorly calibrated. Models routinely express high confidence on questions they answer incorrectly. This overconfidence undermines trust in LLM outputs, particularly for applications where users must decide whether to rely on model predictions.

Standard prompting provides no mechanism for uncertainty to surface during generation. When prompted to answer directly, the model produces a response without intermediate reasoning that might reveal epistemic uncertainty. Confidence scores, if requested, are generated without access to information about the model's uncertainty during answer generation.

\subsubsection{Research Gap}

Prior work has examined chain-of-thought prompting for accuracy improvement (Wei et al., 2022; Kojima et al., 2022) and confidence elicitation methods separately (Lin et al., 2022; Xiong et al., 2023; Tian et al., 2023). However, no systematic study has isolated how the combination of CoT and confidence verbalization might improve calibration, nor whether any improvement stems from meaningful uncertainty signals versus artifacts such as increased token count.

\subsubsection{Hypothesized Mechanism}

This study decomposes the hypothesized mechanism into four testable steps:

\begin{enumerate}
\item CoT prompting produces multi-step reasoning chains.
\item Reasoning chains contain epistemic hedging markers (e.g., ''may,'' ''could,'' ''however'').
\item Due to autoregressive generation, markers are present in-context when confidence is generated.
\item The model incorporates these signals, producing lower confidence when more hedging is present.
\end{enumerate}

\subsubsection{Contributions}

This paper provides:
\begin{enumerate}
\item Empirical validation of each step in the mechanism chain linking CoT prompting to hedging-confidence correlation.
\item Quantification of the hedging-confidence relationship (Spearman r = -0.315) on TruthfulQA.
\item Documentation of structural prerequisites for prompting-based calibration via mechanism decomposition.
\end{enumerate}

The primary ECE comparison across prompting conditions was not executed with real API calls; thus, end-to-end calibration improvement is not demonstrated.
